\documentclass[conference]{IEEEtran}
\IEEEoverridecommandlockouts
\usepackage{cite}
\usepackage{amsmath,amssymb,amsfonts}
\usepackage{graphicx}
\usepackage{textcomp}
\usepackage{xcolor}
\usepackage{hyperref}
\usepackage{tabularx}
\usepackage{booktabs}

\begin{document}

\title{VibeGuardBench: A Deterministic Benchmark for \\ Repair-and-Regression in AI-Generated Web Applications}

\author{\IEEEauthorblockN{Sayed Ali Haider Shah}
\IEEEauthorblockA{Independent Researcher \\ \texttt{github.com/iamcipherdev} \\ October 2026}}

\maketitle

\begin{abstract}
Prompt-driven AI application builders can produce convincing full-stack web
applications from a single specification. Existing benchmarks score that first
act: what the builder generates. What happens next --- when deterministic
engineering feedback identifies concrete defects and the builder is asked to fix
them --- is unmeasured. VibeGuardBench is an open benchmark for that second act.
Three frozen application specifications (appointment booking, multi-role
dashboard, ordering/inventory) are generated by commercial AI builders; a
deterministic 39-check oracle suite (build, types, lint, functional, authorization,
security, accessibility, persistence) measures each application; machine-generated
feedback is delivered per failing check; and the \emph{entire} suite re-executes
after every repair, for up to three iterations. The benchmark reports repair
success rate, regression rate among previously passing checks, and net quality
change, with raw counts and a preregistered analysis plan. A harness-validation
pilot (one application, excluded from research claims) exercised the full loop and
surfaced fifteen harness defects before data collection. The benchmark, oracles,
schemas, and analysis scripts are open source.
\end{abstract}

\begin{IEEEkeywords}
AI-generated code, vibe coding, benchmarks, automated program repair, regression testing, empirical software engineering
\end{IEEEkeywords}

\section{Introduction}
The evaluation of AI programming tools has a blind spot, and it sits exactly where
practitioners live. Benchmarks such as WebGen-Bench~\cite{lu2025webgen},
FullStackBench~\cite{byte2025fullstack}, and AppWorld~\cite{trivedi2024appworld}
score what an AI builder produces from a specification. SWE-bench and its
successors~\cite{jimenez2024swebench} score patches against unit tests in human
repositories. None of them measures the ordinary, expensive middle of real
software work: the generated application fails objective checks, someone hands the
failures back, and the builder tries again. Does the fix work? And what else broke
in the process?

Classical automated program repair learned this lesson the hard way. Patches that
satisfy a test oracle frequently break behavior elsewhere --- documented as test
overfitting~\cite{smith2015cure}, patch overfitting~\cite{petke2024patch}, and
preservation violations~\cite{ismayilzada2023poracle}. Whether modern AI
application builders exhibit the same behavior when repairing their own output, at
the level of a whole application with functional, security, accessibility, build,
and persistence concerns, has not been measured. VibeGuardBench exists to measure
it.

The benchmark's research question is deliberately narrow: \textit{when
deterministic engineering feedback identifies failures in an AI-generated
full-stack application, can the same AI builder repair those failures without
introducing regressions in previously passing functionality?} Everything in the
design serves that question. Three commitments define it. First,
\emph{deterministic measurement}: every check produces the same outcome given the
same artifact; no LLM judges anything. Second, the \emph{whole-suite rule}: after
every repair, all checks re-execute, because regressions are only visible to
checks that run. Third, \emph{pre-freezing}: specifications, prompts, oracles,
metric definitions, and the analysis plan were committed publicly before any
experimental data existed.

This paper describes the benchmark as a tool: the tasks, the oracle suite, the
repair protocol, the metrics, the implementation, and the pilot that validated the
harness. A companion protocol document~\cite{vgbprotocol} records the frozen study
design; this paper is about the instrument itself.

\section{Related Work}
Benchmarks for AI-generated software cluster around initial generation.
WebGen-Bench~\cite{lu2025webgen} scores generated websites with automated tests;
FullStackBench~\cite{byte2025fullstack} spans full-stack task domains;
AppWorld~\cite{trivedi2024appworld} verifies agent programs against simulated
applications; DevBench~\cite{li2024devbench} covers the development lifecycle.
SWE-bench~\cite{jimenez2024swebench} and its agents~\cite{yang2024sweagent,
xia2024agentless, zhang2024autocoderover} grade patches to human repositories
against fail-to-pass tests --- close in spirit to repair, but the subject is a
patch to existing code, not a builder repairing its own generated application,
and the oracle is a unit-test subset rather than a whole-application suite.

The repair literature proper is older and sobering. Smith et al.~\cite{smith2015cure}
showed that most ``repaired'' programs in a large APR study were not actually
correct; subsequent work formalized overfitting as the central pathology of
test-driven repair~\cite{petke2024patch, lou2024apr, ismayilzada2023poracle}.
Self-repair loops for LLM code (Self-Refine~\cite{madaan2023selfrefine},
Self-Correct~\cite{huang2023selfcorrect}, Reflexion~\cite{shinn2023reflexion})
report improved pass rates but rarely re-run a full behavioral suite, so
regressions are structurally undercounted --- the exact flaw the whole-suite rule
is designed to prevent.

On security evaluation, Pearce et al.~\cite{pearce2022asleep, pearce2023zeroshot}
and Perry et al.~\cite{perry2023insecure} established that LLM-generated code
carries distinctive vulnerability patterns; VibeGuardBench folds a frozen static
ruleset and probe-based checks into the suite rather than treating security as a
separate study. A recent survey of vibe coding~\cite{ge2025vibe} identifies
reliability and security of generated products as open empirical problems; this
benchmark is a direct instrument for both.

\section{Benchmark Design}
The design is a factorial grid with a longitudinal component at the check level:
3 application tasks $\times$ 2 commercial builders $\times$ 3 independent
repetitions, yielding 18 initial applications (iteration 0), each eligible for up
to 3 repair iterations. The statistical unit is the check observation
(application $\times$ iteration $\times$ check): 39 checks in the catalog,
roughly 20--25 eligible per suite execution, on the order of 1,500--2,000
observations across 72 executions.

\paragraph{Application tasks.} Three bounded specifications, frozen in the
repository: APP-01, an appointment booking system (two roles, CRUD, search, admin
panel); APP-02, a multi-role SaaS dashboard (three roles, project--task
hierarchy, role-scoped API); APP-03, an ordering/inventory application (order
placement with stock decrement, manager-only inventory administration). Each
specification carries a binding \emph{interface contract} --- exact HTTP routes,
request/response JSON shapes, seed accounts, DOM \texttt{data-testid} attributes
--- which is what makes deterministic cross-implementation testing possible.
Builders keep freedom over stack, architecture, and visual design; contract
deviations fail the affected checks, and that failure is a measured outcome, not
a setup error.

\paragraph{Repair loop.} After each suite execution, a feedback generator renders
one fixed-template message per failing check --- deterministic text with observed
versus expected values and a specification reference --- preceded by a preamble
instructing repair without changing unrelated behavior. The operator delivers the
feedback verbatim to the same builder session, exports the result, and the runner
re-executes the entire suite. Iteration 0 is archived once and never overwritten;
repairs stop early only when nothing is failing.

\section{The Oracle Suite}
Thirty-nine checks across eight categories (Table~\ref{tab:oracles}), each with a
frozen identifier from the check catalog. The runner installs, builds, launches,
and probes the application automatically, recording a schema-validated run record
with raw tool output per check.

\begin{table}[tbp]
\centering
\caption{Oracle categories (39 checks). Runtime categories record \emph{not
applicable} when the application cannot start (cascade rule).}
\label{tab:oracles}
\setlength{\tabcolsep}{4pt}
\small
\begin{tabularx}{\columnwidth}{llX}
\toprule
Category & \# & Method \\
\midrule
BUILD & 5 & fresh install; production build; start on \$PORT; reachability; log hygiene \\
TYPE & 1 & \texttt{tsc --noEmit} \\
LINT & 2 & ESLint, frozen config; zero errors / zero warnings \\
FUNC & 11 & Playwright (Chromium): login, logout, create, list, edit, delete, validation, search, empty state, not-found, API rejection \\
AUTHZ & 6 & role$\times$route API probes; cross-owner writes; UI denial; logout invalidation \\
SEC & 6 & frozen Semgrep ruleset; secret scan; \texttt{npm audit}; debug probes; cookie/token storage; credential storage \\
A11Y & 4 & axe-core (WCAG 2.1 A/AA subset); lang/title \\
DB & 4 & create/update/delete survive process restart; contract schema shapes \\
\bottomrule
\end{tabularx}
\end{table}

Three interpretation rules matter. \emph{Cascade:} unstartable applications record
runtime checks as not applicable; the install/build failures carry the signal.
\emph{Flakiness policy:} environment-class failures are retried once and logged;
test-content failures never are. \emph{Security caution:} scanner findings are
candidate findings, never reported as confirmed vulnerabilities; every distinct
finding is manually inspected, with a false-positive sensitivity analysis in the
frozen plan. No CVSS scores, no exploitability claims from scanner output alone.

\section{Metrics}
Per application $a$ and transition $k \to k{+}1$, let $T(a,k)$ be failing checks
and $P(a,k)$ passing checks at iteration $k$ (environment-error and
not-applicable checks excluded):
\begin{align}
\mathrm{RSR}(a,k) &= \frac{|\{c \in T(a,k): X_{a,k+1}(c)=\mathrm{pass}\}|}{|T(a,k)|}, \\
\mathrm{RR}(a,k) &= \frac{|\{c \in P(a,k): X_{a,k+1}(c)=\mathrm{fail}\}|}{|P(a,k)|}, \\
\mathrm{NQC}(a,k) &= |\mathrm{repaired}| - |\mathrm{regressed}|.
\end{align}
RSR (repair success rate), RR (regression rate), and NQC (net quality change)
are the primary estimands. Raw counts accompany every rate; small denominators
are flagged; application-level clustering handles the non-independence of checks
within an application. The analysis plan is estimation-first with cluster
bootstraps, pre-declared hypothesis families with Benjamini--Hochberg correction,
and effect sizes as risk differences --- frozen before data collection.

\section{Implementation}
The repository ships the complete instrument: frozen specifications and prompts;
the 39-check catalog with per-check feedback templates; oracle implementations
(Playwright suites, a runtime probe engine, the Semgrep ruleset, ESLint config);
\texttt{run\_oracles.py} (full-suite executor), \texttt{make\_feedback.py}
(deterministic feedback), \texttt{ledger.py} (raw $\to$ processed);
\texttt{metrics.py} and \texttt{plots.py} (analysis); JSON schemas for run
records and check results; and an operator runbook for the builder-facing steps.
Every suite execution writes environment versions, raw tool outputs, operator
interventions, and retries into the run record. Every number in the results
regenerates from raw data by script; none are transcribed by hand.

\section{Pilot Validation}
Before any study data, the harness was validated end to end on a pilot
application: APP-01 generated by a non-study LLM system (explicitly excluded from
research claims), three full suite executions, two real repair loops. The
corrected trajectory: 19 passing / 16 failing at iteration 0; a repair that
changed code yet fixed 0 of 16 failures while regressing 1 of 19 passing checks
(the builder moved the session token into a cookie without the required
\texttt{HttpOnly}/\texttt{SameSite} flags); and a correct repair of that
regression at iteration 2. The pilot's value was not the trajectory but what it
exposed: fifteen harness defects --- a lint configuration producing spurious
errors, an over-strict health check misreading redirects, orphaned server
processes, an invalid scanner configuration, silent pass-through of
dependency-audit failures, cross-iteration state pollution, and nine more ---
each fixed and logged before data collection. A benchmark that has not survived
contact with a real repair loop is a specification, not an instrument; the pilot
is what makes VibeGuardBench the latter.

\section{Limitations}
Two builders, three tasks, one operator bound generalization; conclusions concern
the studied systems under the studied contracts. The prescriptive interface
contract trades ecological representativeness for measurement validity. Builders
are compared as systems, not models, since internals are undisclosed. Three
repair iterations bound trajectory claims. Platform versions will drift during
data collection; run records are the only reliable account of which version
produced which artifact.

\section{Conclusion}
Initial-generation benchmarks measure the first act of AI-assisted development.
VibeGuardBench instruments the second: repair under deterministic feedback, with
regressions counted by construction rather than by accident. The benchmark is
open, the oracles are frozen, the analysis is preregistered, and the harness has
survived a real repair loop. The study it enables --- whether AI application
builders can fix their own output without breaking what already works --- is,
to our knowledge, unmeasured anywhere else.

\begin{thebibliography}{00}
\bibitem{ge2025vibe} Y.~Ge et al., ``A survey of vibe coding with large language models,'' \textit{arXiv:2510.12399}, 2025.
\bibitem{lu2025webgen} J.~Lu et al., ``WebGen-Bench: Evaluating LLM-generated web applications,'' 2025.
\bibitem{byte2025fullstack} ByteDance, ``FullStackBench,'' 2025.
\bibitem{trivedi2024appworld} H.~Trivedi et al., ``AppWorld: A controllable world of apps,'' 2024.
\bibitem{jimenez2024swebench} C.~Jimenez et al., ``SWE-bench: Can language models resolve real-world GitHub issues?'' in \textit{Proc.\ ICLR}, 2024.
\bibitem{yang2024sweagent} J.~Yang et al., ``SWE-agent: Agent-computer interfaces enable automated software engineering,'' 2024.
\bibitem{xia2024agentless} C.~Xia et al., ``Agentless: Demystifying LLM-based software engineering agents,'' 2024.
\bibitem{zhang2024autocoderover} Y.~Zhang et al., ``AutoCodeRover: Autonomous program improvement,'' 2024.
\bibitem{smith2015cure} E.~Smith et al., ``Is the cure worse than the disease?'' in \textit{Proc.\ SOSP}, 2015.
\bibitem{ismayilzada2023poracle} S.~Ismayilzada et al., ``Preservation violations in APR,'' 2023.
\bibitem{petke2024patch} J.~Petke et al., ``Patch overfitting in automated program repair,'' 2024.
\bibitem{lou2024apr} Y.~Lou et al., ``Automated program repair: A survey,'' 2024.
\bibitem{madaan2023selfrefine} A.~Madaan et al., ``Self-refine: Iterative refinement with self-feedback,'' 2023.
\bibitem{huang2023selfcorrect} J.~Huang et al., ``Large language models cannot self-correct reasoning yet,'' 2023.
\bibitem{shinn2023reflexion} N.~Shinn et al., ``Reflexion: Language agents with verbal reinforcement learning,'' 2023.
\bibitem{pearce2022asleep} H.~Pearce et al., ``Asleep at the keyboard?'' in \textit{Proc.\ CCS}, 2022.
\bibitem{pearce2023zeroshot} H.~Pearce et al., ``Zero-shot vulnerability repair,'' 2023.
\bibitem{perry2023insecure} N.~Perry et al., ``Do users write more insecure code with AI assistants?'' 2023.
\bibitem{vgbprotocol} S.~A.~H.~Shah, ``VibeGuardBench research protocol --- study in progress,'' \textit{github.com/iamcipherdev/VibeGuardBench}, v1.0.2, Sept.\ 2026.
\end{thebibliography}

\end{document}
